\chapter{Bunionectomy}

\section*{Introduction \& Clinical Indications}
A bunionectomy is a surgical procedure designed to correct hallux valgus, a complex progressive deformity characterized by the lateral deviation of the great toe (hallux) at the first metatarsophalangeal (MTP) joint. This condition is typically associated with a painful bony prominence, known as a bunion, on the medial aspect of the foot, resulting from the medial deviation of the first metatarsal. The deformity can lead to significant pain, inflammation, difficulty with footwear, and impaired ambulation. When conservative management strategies prove insufficient to alleviate symptoms or halt progression, a bunionectomy aims to restore the anatomical alignment of the first MTP joint by resecting the bony exostosis and, crucially, performing osteotomies and soft tissue balancing to realign the first metatarsal and proximal phalanx. This intervention is clinically significant as it can dramatically improve patient quality of life by reducing pain, enhancing mobility, and allowing for comfortable shoe wear.

\begin{itemize}
  \item Hallux valgus correction
  \item Hallux valgus repair
  \item Bunion surgery
  \item First metatarsal osteotomy
  \item Chevron osteotomy
  \item Scarf osteotomy
  \item Akin osteotomy
  \item Lapidus procedure (first tarsometatarsal arthrodesis)
\end{itemize}

The fundamental principle of bunionectomy is to restore the normal anatomical alignment and biomechanical function of the first metatarsophalangeal joint. The rationale for surgical intervention stems from the understanding that hallux valgus is a multiplanar deformity involving both bony malalignment and soft tissue imbalance. Simply removing the bony prominence (exostectomy) without addressing the underlying skeletal and soft tissue pathology often leads to recurrence.

The operative concept involves several key components:
\begin{itemize}
  \item Resection of the medial bony exostosis: This removes the painful bump and allows for proper soft tissue rebalancing.
  \item Osteotomy of the first metatarsal: This is the cornerstone of most modern bunionectomies. By cutting and repositioning the first metatarsal bone, the surgeon can correct the increased intermetatarsal angle (the angle between the first and second metatarsals) and the hallux valgus angle (the angle between the first metatarsal and the proximal phalanx). Common osteotomy types include distal osteotomies (e.g., Chevron, Reverdin) for mild-to-moderate deformities, and proximal osteotomies (e.g., Ludloff, Scarf, basal wedge) or fusion (Lapidus procedure) for more severe deformities or hypermobility.
  \item Soft tissue balancing: This involves releasing contracted lateral structures, such as the adductor hallucis tendon and the lateral joint capsule, and tightening the medial joint capsule (medial capsular plication). This rebalances the forces acting on the great toe, helping to maintain the corrected alignment.
  \item Akin osteotomy: Often performed in conjunction with a metatarsal osteotomy, this is a medial closing wedge osteotomy of the proximal phalanx to correct any residual hallux valgus interphalangeus.
\item Fixation: The osteotomized bone segments are stabilized with screws, pins, or wires to allow for proper healing in the corrected position.
\end{itemize}
The technical goals are to achieve a straight, pain-free great toe with stable correction, preserved joint motion, and the ability to wear normal footwear without discomfort or recurrence.

The surgical correction of hallux valgus has a rich history spanning over a century, marked by a progressive understanding of the deformity's complexity and the development of increasingly sophisticated techniques. Early attempts in the late 19th and early 20th centuries, such as those described by Reverdin in 1881, primarily focused on simple exostectomy (removal of the bony prominence) and sometimes a medial closing wedge osteotomy of the metatarsal head. While these provided initial relief, high recurrence rates highlighted the need to address the underlying metatarsal malalignment and soft tissue imbalance.

A significant advancement came with the McBride procedure in 1928, which involved lateral soft tissue release and adductor hallucis transfer, sometimes combined with exostectomy. This recognized the importance of soft tissue balancing. The mid-20th century saw the introduction of various metatarsal osteotomies. The Chevron osteotomy, popularized by Austin and Leventen in 1962, became a widely adopted distal osteotomy for mild to moderate deformities due to its stability and predictable outcomes. The Scarf osteotomy, developed by Barouk in 1985, offered a versatile, long Z-cut osteotomy of the metatarsal shaft, allowing for significant correction in multiple planes and stable fixation, suitable for a broader range of deformities. More recently, the Lapidus procedure, a fusion of the first tarsometatarsal joint, has gained favor for severe deformities, hypermobility of the first ray, or recurrent cases, addressing the root cause of proximal metatarsal instability. The evolution continues with minimally invasive techniques and improved fixation methods, all aiming for better patient outcomes and reduced recovery times.

A bunionectomy is indicated for patients experiencing significant symptoms and functional limitations due to hallux valgus, particularly when conservative measures have failed. Specific clinical indications include:

\begin{itemize}
  \item Persistent pain over the medial bunion prominence, especially exacerbated by walking or pressure from footwear, which significantly impacts daily activities.
  \item Difficulty finding comfortable or appropriate shoes due to the size and shape of the bunion and the deviation of the great toe.
  \item Progressive deformity that interferes with ambulation, balance, or leads to secondary deformities such as hammertoes of the lesser toes or crossover toe.
  \item Chronic inflammation, redness, or skin irritation (bursitis) over the bunion, unresponsive to protective padding.
  \item Stiffness or limited range of motion of the first metatarsophalangeal joint, contributing to pain and functional impairment.
  \item Neurological symptoms such as numbness or tingling in the great toe due to nerve compression from the deformity.
  \item Radiographic evidence of significant hallux valgus angle (HVA > 20 degrees) and/or intermetatarsal angle (IMA > 9 degrees), especially if progressive.
  \item Failure of a comprehensive trial of conservative management, including wide-toed shoes, orthotics, toe spacers, padding, and anti-inflammatory medications, typically over a period of 3-6 months.
\end{itemize}

A bunionectomy is considered the primary definitive treatment for symptomatic hallux valgus, but it is positioned as a secondary intervention in the overall management pathway. This means that surgical correction is typically reserved for patients who have failed an adequate trial of conservative, non-operative measures. For individuals with mild or asymptomatic hallux valgus, or those whose symptoms are effectively managed, surgery is generally not recommended.

Conservative care, including appropriate footwear modifications (wide toe box, low heel), orthotic devices, toe spacers, padding, and activity modification, is the first-line approach. If these measures successfully alleviate pain and prevent progression, surgery may be avoided indefinitely. However, for patients with persistent pain, significant functional impairment, or progressive deformity despite conservative efforts, a bunionectomy becomes the primary and most effective treatment to correct the underlying anatomical pathology. It is not a salvage procedure but rather a reconstructive one aimed at restoring normal foot mechanics and improving quality of life. The decision to proceed with surgery is highly individualized, balancing the severity of symptoms, the degree of deformity, patient expectations, and potential risks.

\section*{Relevant Surgical Anatomy}
The primary anatomical focus of a bunionectomy is the first metatarsophalangeal (MTP) joint, which connects the first metatarsal bone of the foot to the proximal phalanx of the great toe. Key bony structures include the head of the first metatarsal, the base of the proximal phalanx, and the two sesamoid bones embedded within the plantar plate beneath the metatarsal head. In hallux valgus, the first metatarsal deviates medially, while the great toe deviates laterally, leading to subluxation of the sesamoids and a prominent medial eminence on the metatarsal head.

Surrounding soft tissues are critical to the deformity and its correction. The joint capsule is typically stretched and attenuated medially, while it becomes contracted and thickened laterally. The adductor hallucis muscle, which inserts into the lateral aspect of the proximal phalanx and the lateral sesamoid, contributes to the lateral deviation of the toe and the medial deviation of the metatarsal. Conversely, the abductor hallucis muscle, inserting medially, becomes functionally weakened or displaced. The medial and lateral collateral ligaments of the MTP joint are also affected, with the medial ligament often lax and the lateral ligament contracted. Neurovascular structures requiring careful identification and protection during surgery include the medial dorsal cutaneous nerve, a branch of the superficial peroneal nerve, which supplies sensation to the dorsomedial aspect of the great toe, and the medial plantar nerve and digital arteries, which run along the plantar aspect. Surgical landmarks include the palpable medial eminence and the dorsal aspect of the first MTP joint capsule.

\section*{Preoperative Workup \& Preparation}
Thorough preoperative workup and preparation are essential to optimize outcomes and minimize risks for bunionectomy. This begins with a detailed medical history and physical examination, focusing on the patient's overall health, comorbidities, and specific foot pathology.

\begin{itemize}
  \item \textbf{Medical Evaluation:} A comprehensive medical evaluation is performed to assess cardiac, pulmonary, and renal function, especially for older patients or those with chronic diseases. This may include an electrocardiogram (ECG), chest X-ray, and consultation with an internist or anesthesiologist.
  \item \textbf{Laboratory Tests:} Routine blood tests typically include a complete blood count (CBC), basic metabolic panel (BMP), and coagulation studies (PT/INR, PTT). For diabetic patients, HbA1c levels are checked to ensure adequate glycemic control.
  \item \textbf{Imaging Studies:} Weight-bearing radiographs of the foot (anteroposterior, lateral, and oblique views) are mandatory to accurately assess the degree of hallux valgus angle, intermetatarsal angle, sesamoid position, and joint congruity, which guides surgical planning.
  \item \textbf{Medication Adjustment:} Patients on anticoagulants or antiplatelet agents (e.g., aspirin, clopidogrel, warfarin) will receive specific instructions from their surgeon or prescribing physician regarding temporary cessation or bridging therapy to minimize bleeding risk.
  \item \textbf{NPO Status:} Patients are instructed to be NPO (nil per os) for a specified period (typically 6-8 hours) before surgery to prevent aspiration during anesthesia.
  \item \textbf{Prophylactic Antibiotics:} Intravenous prophylactic antibiotics (e.g., cefazolin) are administered within one hour prior to incision to reduce the risk of surgical site infection.
  \item \textbf{Informed Consent:} A detailed discussion with the patient covers the surgical procedure, expected outcomes, potential risks, and alternative treatments, ensuring fully informed consent is obtained.
  \item \textbf{Foot Preparation:} Any active skin infections, fungal infections, or open wounds on the foot must be treated and resolved prior to surgery. Patients are advised on postoperative care, including activity restrictions, pain management, and the need for assistance at home.
\end{itemize}

Careful consideration of contraindications and precautions is crucial for patient safety and successful outcomes in bunionectomy.

\textbf{Absolute Contraindications:}
\begin{itemize}
  \item Active infection in the foot or ankle, including cellulitis, osteomyelitis, or an infected ulcer, as this significantly increases the risk of postoperative infection.
  \item Severe peripheral vascular disease (PVD) with critical limb ischemia, which compromises wound healing and increases the risk of nonunion or amputation.
  \item Uncontrolled Charcot neuroarthropathy, leading to bone fragility and high risk of nonunion or further deformity.
  \item Unrealistic patient expectations regarding cosmetic outcome or complete pain relief, which can lead to dissatisfaction despite a technically successful surgery.
  \item Inability or unwillingness of the patient to comply with postoperative weight-bearing restrictions and rehabilitation protocols.
\end{itemize}

\textbf{Relative Contraindications \& Precautions:}
\begin{itemize}
  \item Uncontrolled diabetes mellitus: Requires meticulous glycemic control preoperatively and careful monitoring postoperatively due to increased risk of infection, poor wound healing, and neuropathy.
  \item Significant neuropathy (e.g., diabetic neuropathy): May mask pain, leading to excessive weight-bearing and potential complications like stress fractures or nonunion.
  \item Severe osteoporosis: Increases the risk of intraoperative fracture and challenges with internal fixation, potentially requiring alternative fixation methods or prolonged non-weight-bearing.
  \item Active smoking: Significantly impairs bone healing and increases the risk of nonunion, infection, and wound complications. Smoking cessation is strongly advised preoperatively.
  \item Mild or asymptomatic deformity: Surgery is generally not indicated for cosmetic reasons alone or for deformities that do not cause significant pain or functional impairment.
  \item Severe systemic disease: Conditions that make anesthesia unsafe or significantly compromise recovery (e.g., severe cardiac or pulmonary disease) require careful risk-benefit assessment.
  \item Very elderly or frail patients: Increased risk of complications and slower recovery may warrant a more conservative approach.
\end{itemize}
Precautions include meticulous soft tissue handling, careful hemostasis, and vigilant monitoring of neurovascular status, particularly in patients with compromised circulation or sensation.

\section*{Surgical Technique \& Procedure}
A bunionectomy is typically performed under regional anesthesia (e.g., ankle block or popliteal block) combined with sedation, or general anesthesia, often as an outpatient procedure. The patient is positioned supine, and a thigh tourniquet is usually applied to create a bloodless field, enhancing surgical visibility. Prophylactic antibiotics are administered intravenously prior to incision.

The general steps for a common bunionectomy (e.g., Chevron or Scarf osteotomy) include:
\begin{itemize}
  \item \textbf{Incision:} A longitudinal incision is made on the medial aspect of the first metatarsophalangeal joint, carefully avoiding the medial dorsal cutaneous nerve. The incision extends from the mid-shaft of the first metatarsal to the mid-shaft of the proximal phalanx.
  \item \textbf{Capsulotomy and Exostectomy:} The joint capsule is incised, typically in a medial longitudinal or L-shaped fashion. The medial bony prominence (exostosis) on the first metatarsal head is then exposed and resected using an oscillating saw, ensuring a smooth contour.
  \item \textbf{Lateral Release:} A critical step for correcting the hallux valgus angle involves releasing the contracted lateral structures. This typically includes a release of the adductor hallucis tendon from its insertion on the proximal phalanx and lateral sesamoid, and a lateral capsulotomy of the first MTP joint. This allows the great toe to be realigned.
  \item \textbf{Osteotomy:} The specific type of osteotomy depends on the severity and location of the deformity.
    \begin{itemize}
      \item \textit{Chevron Osteotomy:} A V-shaped cut is made in the distal metatarsal head, allowing the capital fragment to be shifted laterally to reduce the intermetatarsal angle.
      \item \textit{Scarf Osteotomy:} A Z-shaped cut is made along the shaft of the first metatarsal, allowing for translation, rotation, and shortening/lengthening of the bone, providing greater correction potential.
      \item \textit{Basal/Proximal Osteotomy:} A wedge or crescentic cut is made near the base of the first metatarsal for more severe deformities.
      \item \textit{Akin Osteotomy:} A medial closing wedge osteotomy of the proximal phalanx may be performed to correct any residual hallux valgus interphalangeus.
    \end{itemize}
  \item \textbf{Fixation:} Once the bone segments are realigned, they are stabilized with internal fixation, typically using one or two small screws, K-wires, or absorbable pins, to ensure proper healing.
  \item \textbf{Medial Capsular Plication:} The medial joint capsule is tightened (plicated) to reinforce the medial side of the joint and help maintain the corrected alignment.
  \item \textbf{Wound Closure:} The tourniquet is deflated, hemostasis is achieved, and the wound is closed in layers. A sterile dressing and a compressive bandage are applied, often followed by a rigid-soled surgical shoe or boot. Drains are rarely used.
\end{itemize}
The procedure duration varies but typically ranges from 45 to 90 minutes.

\section*{Complications \& Risks}
While bunionectomy is generally safe, like any surgical procedure, it carries potential complications and risks. These can be broadly categorized into general surgical risks and procedure-specific risks.

\textbf{General Surgical Complications:}
\begin{itemize}
  \item \textbf{Anesthesia-related risks:} Reactions to anesthetic agents, respiratory depression, nausea, vomiting, or, rarely, more severe cardiovascular events.
  \item \textbf{Bleeding/Hematoma:} Excessive bleeding during or after surgery, potentially requiring drainage.
  \item \textbf{Infection:} Surgical site infection (superficial or deep osteomyelitis), which may require antibiotics or further surgery.
  \item \textbf{Deep Vein Thrombosis (DVT) / Pulmonary Embolism (PE):} Blood clot formation, though rare in foot surgery, can occur, especially with prolonged immobilization.
  \item \textbf{Wound healing problems:} Delayed healing, dehiscence, or scarring.
\end{itemize}

\textbf{Procedure-Specific Complications \& Risks:}
\begin{itemize}
  \item \textbf{Recurrence of the deformity:} The bunion may return over time, especially if the underlying biomechanical factors are not fully addressed or if postoperative care is not followed.
  \item \textbf{Undercorrection or Overcorrection:} The great toe may not be adequately straightened (undercorrection) or may be straightened too much (hallux varus), leading to new problems.
  \item \textbf{Nonunion or Malunion of the osteotomy:} The bone cut may fail to heal (nonunion) or heal in an undesirable position (malunion), potentially requiring revision surgery.
  \item \textbf{Avascular Necrosis (AVN) of the metatarsal head:} Particularly a risk with distal osteotomies (e.g., Chevron), where blood supply to the metatarsal head can be compromised, leading to bone death and joint collapse.
  \item \textbf{Nerve injury:} Damage to the medial dorsal cutaneous nerve can cause numbness, tingling, or painful neuroma formation on the dorsomedial aspect of the great toe.
  \item \textbf{Stiffness of the first MTP joint:} Reduced range of motion, which can affect walking and shoe wear.
  \item \textbf{Transfer Metatarsalgia:} Pain under the lesser metatarsal heads (typically the second or third) due to altered weight-bearing patterns after correction of the first ray.
  \item \textbf{Hardware irritation:} Screws or pins used for fixation may become prominent or cause discomfort, sometimes requiring removal after bone healing.
  \item \textbf{Complex Regional Pain Syndrome (CRPS):} A rare but severe chronic pain condition affecting the limb.
\end{itemize}

\section*{Postoperative Care \& Recovery}
The postoperative recovery timeline for a bunionectomy is a gradual process, typically extending over several months, with distinct phases of healing and rehabilitation.

Immediately after surgery, the patient is transferred to the Post-Anesthesia Care Unit (PACU) for close monitoring of vital signs, pain levels, and neurovascular status of the foot. The foot will be bandaged and elevated to minimize swelling. Pain management is initiated with oral or intravenous analgesics. Within the first 24-48 hours, the patient is typically discharged home with instructions for strict elevation, icing, and non-weight-bearing or partial weight-bearing in a specialized surgical shoe or boot, depending on the surgeon's protocol and the type of osteotomy performed. Crutches or a walker are often required for mobility. Pain and swelling are expected, managed with prescribed medications and continued elevation.

Over the first 1-2 weeks, wound care involves keeping the dressing clean and dry. The initial bulky dressing is usually changed by the surgeon or nurse, and sutures or staples are typically removed around 10-14 days post-op. During this period, weight-bearing restrictions remain crucial to allow the osteotomy to begin healing. Patients are advised to avoid prolonged standing or walking. Gradual progression to full weight-bearing in the surgical shoe or boot usually occurs between 4-8 weeks, depending on radiographic evidence of bone healing. Physical therapy may be initiated to restore range of motion and strength. Long-term recovery, involving the resolution of residual swelling and stiffness, can take 3-6 months, with some patients experiencing mild swelling for up to a year. Return to normal activities, including sports, is typically permitted after 3-6 months, once bone healing is complete and strength is regained.

During a bunionectomy, the surgeon typically observes several characteristic findings that confirm the diagnosis and guide the corrective steps. The most prominent finding is the hypertrophied medial eminence (bony exostosis) on the head of the first metatarsal. The first metatarsophalangeal joint capsule is often found to be attenuated and redundant on the medial side, while being tight and contracted on the lateral side. The great toe itself is visibly deviated laterally, and the sesamoid bones, normally positioned directly beneath the metatarsal head, are often subluxed laterally. The adductor hallucis tendon is typically taut and contributes to the lateral pull on the great toe.

Pathological examination of the resected medial eminence, if sent for analysis, typically reveals mature cortical and cancellous bone, often with evidence of degenerative changes suchosis as cartilage fibrillation, osteophyte formation, and subchondral sclerosis, consistent with osteoarthritis secondary to the chronic mechanical stress of the deformity. There are usually no neoplastic or infectious findings unless specifically suspected preoperatively. The pathology report confirms the benign nature of the bony prominence and provides histological context to the degenerative process occurring within the joint.

A normal and successful postoperative course following a bunionectomy is characterized by a progressive reduction in pain, gradual resolution of swelling, and restoration of functional mobility. Initially, pain is managed with medication, but it should steadily decrease over the first few weeks. Swelling, which can be significant in the immediate postoperative period, will slowly subside, though some residual swelling may persist for several months. The osteotomy site will begin to heal, with radiographic evidence of callus formation typically visible by 4-6 weeks.

Patients should expect to follow a structured rehabilitation program, gradually increasing weight-bearing and activity levels as guided by the surgeon. The great toe should appear straighter, and the painful medial prominence should be absent. By 3-6 months, most patients experience significant improvement in pain, can wear a wider range of shoes comfortably, and can resume most daily activities. The ultimate goal is a stable, pain-free, and well-aligned great toe that allows for comfortable ambulation and improved quality of life, without recurrence of the deformity or transfer of pain to other parts of the foot.

Postoperative care and follow-up are critical for ensuring proper healing, preventing complications, and achieving optimal long-term outcomes after bunionectomy.

\begin{itemize}
  \item \textbf{Wound Care:} Keep the dressing clean and dry. Avoid getting the foot wet until sutures are removed and the incision is well-healed.
  \item \textbf{Elevation and Ice:} Maintain strict elevation of the foot above heart level for the first few days, and apply ice packs to reduce swelling and pain.
  \item \textbf{Weight-Bearing Restrictions:} Adhere strictly to the surgeon's instructions regarding weight-bearing, typically using a surgical shoe or boot and crutches for 4-8 weeks.
  \item \textbf{Pain Management:} Take prescribed pain medications as directed, and transition to over-the-counter analgesics as pain subsides.
  \item \textbf{Follow-up Visits:} Regular follow-up appointments are scheduled, usually at 1-2 weeks for suture/staple removal, and then at 4-6 weeks, 3 months, 6 months, and 1 year post-surgery.
  \item \textbf{Radiographs:} Postoperative X-rays are taken at follow-up visits to monitor osteotomy healing and confirm alignment.
  \item \textbf{Physical Therapy:} May be prescribed to restore range of motion, strength, and proprioception of the great toe and foot, typically starting after initial bone healing.
  \item \textbf{Footwear Advice:} Patients are advised to wear supportive, wide-toed shoes after transitioning out of the surgical boot to prevent recurrence and accommodate any residual swelling.
  \item \textbf{Warning Signs:} Contact the surgeon immediately if experiencing increasing pain, fever, chills, excessive redness or warmth around the incision, pus-like drainage, new numbness or tingling, or inability to move the toes.
\end{itemize}

\section*{Outcomes \& Prognosis}
Hallux valgus is one of the most prevalent forefoot deformities, affecting a significant portion of the adult population worldwide. Its incidence and prevalence vary widely in literature, but generally range from 23\% to 35\% in adults, with a notable increase with age, reaching up to 70\% in individuals over 60 years old. The condition is overwhelmingly more common in women than in men, with a reported female-to-male ratio often cited as high as 10:1 to 15:1. This gender disparity is strongly linked to the prolonged use of narrow, pointed, and high-heeled footwear, which mechanically exacerbates the deformity. Other contributing factors include genetic predisposition (family history), generalized ligamentous laxity, flat feet (pes planus), and certain neuromuscular conditions. While many individuals with hallux valgus remain asymptomatic, a substantial proportion develop pain and functional limitations, making bunionectomy one of the most frequently performed orthopedic procedures of the foot. The overall mortality associated with bunionectomy is exceedingly low, primarily related to general anesthesia risks, while morbidity rates are generally acceptable, with specific complications discussed previously.

The outcomes and prognosis following a bunionectomy are generally favorable, with high rates of patient satisfaction and significant improvement in pain and function. Success rates, defined by pain relief, improved alignment, and ability to wear normal shoes, typically range from 85\% to 95\%. Patients can expect a substantial reduction in pain, correction of the great toe deformity, and enhanced mobility. Functional outcomes include improved gait mechanics and the ability to participate in activities that were previously limited.

However, recurrence of the deformity remains a recognized risk, with rates varying depending on the initial severity, surgical technique, and patient factors, generally ranging from 5\% to 20\%. Other potential long-term issues include persistent stiffness of the MTP joint, transfer metatarsalgia, and hardware-related symptoms. Prognostic factors influencing outcomes include the severity of the preoperative deformity, the presence of arthritis in the MTP joint, patient age, adherence to postoperative rehabilitation, and the surgeon's experience. While there isn't a single landmark clinical trial for bunionectomy akin to those for major joint replacements, numerous studies and systematic reviews consistently demonstrate the efficacy of modern osteotomy-based techniques in achieving lasting correction and patient satisfaction. For instance, studies comparing different osteotomy types (e.g., Chevron vs. Scarf) often show similar functional outcomes but with varying complication profiles and suitability for different deformity severities.

\section*{References}
\begin{enumerate}
  \item MedlinePlus. Bunion removal. U.S. National Library of Medicine; 2024.
  \item Coughlin MJ, Saltzman CL, Anderson RB. Mann's Surgery of the Foot and Ankle. 9th ed. Elsevier; 2013.
  \item American Academy of Orthopaedic Surgeons. Hallux Valgus (Bunion). AAOS Patient Information; 2023.
  \item Barouk LS. Forefoot Reconstruction. 2nd ed. Springer; 2005.
\end{enumerate}

\clearpage
